\section{Reformulation-Enabled Integer L-Shaped Method}
\label{sec:method}

This section develops the proposed method, starting from the lifted reformulation rather than from the natural decomposition baseline.
The method has three linked components: an exact structure-exposing lifted reformulation, a nested scenario--route relaxation--verification decomposition, and joint-resource direction-selective monotonicity cuts.
The notation in this section separates master decisions from scenario operating evaluations.
Let $\mathcal{Q}_\omega(w,z_\omega)$ denote the exact operational feasible set after the battery configuration is fixed, and let $Q_\omega(w,z_\omega)$ denote the corresponding exact operational recourse value.
Let $\mathcal{Q}^R_\omega(w,z_\omega)$ and $Q^R_\omega(w,z_\omega)$ denote the Phase 1 relaxation.
The variable $\theta_\omega$ is the master approximation of $Q_\omega(w,z_\omega)$, and $\mathcal{L}_\omega$ is the current set of valid lower-bounding cuts for scenario $\omega$.

\subsection{Lifted master reformulation}
\label{subsec:lifted-reformulation}

The natural integer L-shaped formulation in Section~\ref{sec:model} keeps only the facility vector $w$ in the master and evaluates each scenario through the battery-and-operation recourse value $\Phi_\omega(w)$.
The lifted reformulation moves the scenario-level battery variables $z_\omega$ into the master so that the master decides the facility layout and the battery configuration together.

For fixed $(w,z_\omega)$, the exact operational recourse value is
\begin{equation}
    Q_\omega(w,z_\omega)
    =
    \min_{y_\omega}
    \left\{
    C^{\mathrm{op}}_\omega(y_\omega)
    :
    y_\omega \in \mathcal{Q}_\omega(w,z_\omega)
    \right\}.
    \label{eq:operational-recourse}
\end{equation}
The scenario value from Section~\ref{sec:model} can then be written as
\begin{equation}
    \Phi_\omega(w)
    =
    \min_{z_\omega}
    \left\{
    C^{\mathrm{bat}}_\omega(z_\omega)
    +
    Q_\omega(w,z_\omega)
    :
    z_\omega \in \mathcal{Z}_\omega
    \right\}.
    \label{eq:lifted-scenario-value}
\end{equation}

The lifted master is
\begin{align}
    \min_{w,z,\theta} \quad
        & C^{\mathrm{fac}}(w)
        + \sum_{\omega \in \Omega} p_\omega
        \left[
            C^{\mathrm{bat}}_\omega(z_\omega) + \theta_\omega
        \right]
        \label{eq:master-obj} \\
    \mathrm{s.t.} \quad
        & z_\omega \in \mathcal{Z}_\omega,
        && \omega \in \Omega,
        \label{eq:master-z-domain} \\
        & \theta_\omega \geq \ell, && \omega \in \Omega,\ \ell \in \mathcal{L}_\omega,
        \label{eq:benders-cuts} \\
        & w_j \in \{0,1\},
        && j \in \mathcal{J},
        \label{eq:master-w-domain} \\
        & z_{\omega r k} \in \{0,1\},
        && \omega \in \Omega,\ r \in \mathcal{R}_\omega,\ k \in \mathcal{K}_r.
        \label{eq:binary-decisions}
\end{align}
Here $\mathcal{Z}_\omega$ is the battery-choice feasible set from Section~\ref{sec:model}.

\begin{proposition}[Exactness of the lifted master reformulation]
\label{prop:equivalence}
The lifted master reformulation is equivalent to the natural integer L-shaped formulation in Section~\ref{sec:model}.
\end{proposition}

\begin{proof}[Proof sketch]
For any feasible solution of the natural formulation, the same $(w,z_\omega)$ pair is feasible in the lifted master, and the exact operational value $Q_\omega(w,z_\omega)$ is unchanged.
Conversely, every feasible lifted solution induces the same $(w,z_\omega)$ configuration and the same objective value in the natural formulation because $z_\omega$ remains scenario dependent and has no cross-scenario nonanticipativity.
\qedhere
\end{proof}

This reformulation only changes the layer at which $z_\omega$ is optimized.
It does not impose cross-scenario nonanticipativity on $z_\omega$ and does not claim novelty from moving recourse variables in general.
Its role is to isolate the remaining operational recourse and expose the facility--battery resource structure used by the Phase 1 decomposition and monotonicity cuts.

\subsection{Two-phase lower bounding and exact verification}
\label{subsec:two-phase-framework}

Given $(w,z_\omega)$, the exact operational subproblem remains difficult because it still contains shared plug-in fast-charger capacity constraints and binary operational decisions.
The two-phase scheme separates lower-bound generation from exact verification.

\subsubsection{Phase 1: route-level LP relaxation}
\label{subsec:phase-one}

Phase 1 relaxes shared plug-in fast-charger capacity coupling and operational integrality in the exact operational subproblem.
As a result, the remaining model decomposes by route and yields smaller LP subproblems whose dual solutions generate lower-bound cuts.
The closer the Phase 1 relaxation is to the Phase 2 feasible region, the tighter these lower bounds can be.
This Phase 1--Phase 2 gap is driven by how binding the relaxed constraints are; in the wireless charging application, the key relaxed coupling is the capacity of existing plug-in fast chargers that serve residual fast-charging demand after wireless charging opportunities and battery choices are fixed.

\begin{proposition}[Phase 1 relaxation and decomposition]
\label{prop:phase-one-lb}
For any fixed master solution $(w,z_\omega)$, the Phase 1 relaxation satisfies
\begin{equation}
    Q^R_\omega(w,z_\omega) \leq Q_\omega(w,z_\omega).
    \label{eq:phase-one-lb}
\end{equation}
Moreover, the Phase 1 relaxation decomposes exactly into independent route-level LP subproblems:
\begin{equation}
    Q^R_\omega(w,z_\omega)
    =
    \sum_{r \in \mathcal{R}_\omega} Q^R_{\omega r}(w,z_{\omega r}).
    \label{eq:phase-one-decomposition}
\end{equation}
\end{proposition}

\begin{proof}[Proof sketch]
Phase 1 removes constraints and integrality requirements from the exact operational subproblem, so its feasible region contains the exact feasible region and its minimization value is no larger.
After the shared plug-in fast-charger capacity coupling is relaxed, the remaining route-specific constraints and objective terms separate by route, which gives the additive decomposition.
\qedhere
\end{proof}

Any LP-dual lower-bounding cut generated from Phase 1 is therefore valid for the exact operational value.
For a violated route-level LP relaxation, a representative cut can be written as
\begin{equation}
    \theta_\omega
    \geq
    \alpha_\omega
    + \sum_{j \in \mathcal{J}} \pi_{\omega j} w_j
    + \sum_{r \in \mathcal{R}_\omega}
      \sum_{k \in \mathcal{K}_r}
      \lambda_{\omega r k} z_{\omega r k},
    \label{eq:phase-one-dual-cut}
\end{equation}
where the coefficients are obtained from the Phase 1 LP dual.

\subsubsection{Phase 2: mixed-integer verification}
\label{subsec:phase-two}

Phase 2 restores the exact operational MIP for fixed $(w,z_\omega)$.
It produces exact feasibility judgments, exact operational costs, and incumbent upper bounds:
\begin{equation}
    UB =
    C^{\mathrm{fac}}(w)
    + \sum_{\omega \in \Omega} p_\omega
    \left[
        C^{\mathrm{bat}}_\omega(z_\omega)
        + Q_\omega(w,z_\omega)
    \right].
    \label{eq:upper-bound}
\end{equation}
Because Phase 2 evaluates the same exact operational subproblem that appears in the lifted formulation, it preserves the exactness of the model.

\subsection{Monotonicity-based strengthened cuts}
\label{subsec:strengthened-cuts}

The full method strengthens both the relaxed and exact separation loops whenever the corresponding monotonicity-based cut is valid for the exact recourse function.
The key structure is the joint partial order induced by wireless facility additions and route-level battery-capacity upgrades.

\begin{assumption}[Monotone feasibility and operating cost]
\label{ass:monotonicity}
For a fixed service-plan scenario, opening additional wireless facilities or selecting weakly larger battery capacities cannot remove any operational schedule that was previously feasible.
Moreover, these changes cannot increase the optimal operational recourse value $Q_\omega(w,z_\omega)$.
\end{assumption}

Assume that the battery options for each route are indexed in nondecreasing capacity order.
Consider a Phase 2 solution $(w^*,z^*_\omega)$ for scenario $\omega$, and let $k^*_{\omega r}$ be the selected battery-capacity option for route $r$ in this solution.
Define
\begin{equation}
    I_\omega(w,z_\omega;w^*,z^*_\omega)
    =
    \sum_{j: w^*_j=0} w_j
    +
    \sum_{r \in \mathcal{R}_\omega}
    \sum_{k \in \mathcal{K}_r: k > k^*_{\omega r}}
    z_{\omega r k}.
    \label{eq:monotone-indicator}
\end{equation}
The indicator is zero when a candidate opens no additional facility and upgrades no route battery relative to $(w^*,z^*_\omega)$.

\begin{proposition}[Strengthened feasibility cut]
\label{prop:strengthened-feasibility}
Suppose Assumption~\ref{ass:monotonicity} holds and the Phase 2 operational subproblem for scenario $\omega$ is infeasible at $(w^*,z^*_\omega)$.
Then
\begin{equation}
    I_\omega(w,z_\omega;w^*,z^*_\omega) \geq 1
    \label{eq:strengthened-feasibility-cut}
\end{equation}
is a valid feasibility cut.
\end{proposition}

\begin{proof}[Proof sketch]
Any solution with zero indicator is dominated by the infeasible configuration and cannot restore feasibility by monotonicity.
\qedhere
\end{proof}

\begin{proposition}[Strengthened optimality cut]
\label{prop:strengthened-optimality}
Suppose Assumption~\ref{ass:monotonicity} holds and the Phase 2 operational subproblem for scenario $\omega$ is feasible at $(w^*,z^*_\omega)$ with value $Q^*_\omega$.
Let $L_\omega$ be any global lower bound satisfying $L_\omega \leq Q_\omega(w,z_\omega)$ for every master-feasible $(w,z_\omega)$, and define $\delta_\omega = Q^*_\omega - L_\omega \geq 0$.
Then
\begin{equation}
    \theta_\omega
    \geq
    Q^*_\omega
    -
    \delta_\omega
    I_\omega(w,z_\omega;w^*,z^*_\omega)
    \label{eq:strengthened-optimality-cut}
\end{equation}
is a valid optimality cut.
\end{proposition}

\begin{proof}[Proof sketch]
For zero indicator, monotonicity gives the bound $Q^*_\omega$.
For positive indicator, the right-hand side relaxes to no more than the global lower bound $L_\omega$.
\qedhere
\end{proof}

These cuts strengthen the lifted master without changing the underlying exact problem.
They are not meant to claim that monotonicity or logic-based Benders cuts are new in general; the contribution is the direction-selective cut structure for the joint facility--battery order.

\subsection{Complete solution procedure}
\label{subsec:complete-procedure}

Algorithm~\ref{alg:full-method} summarizes the full enhanced integer L-shaped method.

\begin{algorithm}[htbp]
\caption{Enhanced integer L-shaped method}
\label{alg:full-method}
\SetKwInOut{Input}{Input}
\SetKwInOut{Output}{Output}
\Input{Scenario data, candidate wireless facilities, candidate battery-capacity options, tolerance $\varepsilon$, time limit $T$}
\Output{Facility layout $w$, scenario battery configurations $z$, incumbent upper bound, and final lower bound}

\BlankLine
Initialize the lifted master with battery-selection constraints and trivial recourse lower bounds\;
\BlankLine
\tcp{Phase 1: route-level LP lower-bounding separation}
\Repeat{no violated Phase 1 cut is found or the time limit is reached}{
    Solve the lifted master\;
    Solve route-level LP relaxations for each scenario at the current $(w,z_\omega)$\;
    Add LP-dual lower-bound cuts and relaxation feasibility cuts\;
    Add monotonicity-based lower-bounding cuts that are valid for exact recourse\;
}
\BlankLine
\tcp{Phase 2: exact mixed-integer verification}
\Repeat{the incumbent gap is at most $\varepsilon$ or the time limit is reached}{
    Solve the lifted master\;
    Solve exact operational MIP subproblems for each scenario at the current $(w,z_\omega)$\;
    Add monotonicity-based feasibility or optimality cuts when violated\;
    Update the incumbent upper bound if all scenarios are feasible\;
}
\Return{Best incumbent and final lower bound}\;
\end{algorithm}

The implementation terminates when the incumbent gap reaches the prescribed tolerance or the time limit is reached.
This stopping rule is distinct from the exact finite-convergence statement below.

\begin{theorem}[Finite convergence]
\label{thm:finite-convergence}
Assume that the lifted master is solved exactly at each iteration, Phase 2 scenario subproblems are solved exactly, and all added feasibility and optimality cuts are valid.
Then the exact Phase 2 separation loop terminates finitely with an optimal solution of the lifted reformulation and hence of the natural integer L-shaped formulation.
\end{theorem}

\begin{proof}[Proof sketch]
The lifted master has finitely many binary facility and battery configurations.
Valid feasibility cuts exclude infeasible configurations.
Valid optimality cuts correct feasible configurations whose exact recourse value is underestimated by the master.
Thus the same unresolved configuration cannot recur indefinitely under exact separation, and Proposition~\ref{prop:equivalence} transfers convergence to the natural integer L-shaped formulation.
\qedhere
\end{proof}
